0% !TeX root = main.tex
\chapter{Introduction}
\label{chap:introduction}

\section{Background and Motivation}
\label{sec:background-motivation}

Real-time speech-to-speech translation is motivated by a straightforward but technically demanding goal: reduce language barriers without forcing users into delayed, text-only, or manually mediated interaction. The societal relevance of multilingual access is well established in educational and public-service contexts, where language inclusion affects participation and access to knowledge \cite{unesco2022inclusion}. At the same time, market reports and recent industrial overviews indicate strong commercial interest in speech translation, especially for cross-border collaboration, customer support, and meeting assistance \cite{mordor2024speech,kudo2024stats}. The combination of social demand and commercial momentum makes real-time translation an important applied research problem rather than a narrow laboratory exercise.

Current production-grade offerings are primarily provided by large cloud vendors through tightly integrated service portfolios such as Google Cloud Translation, Microsoft Translator, and the AWS speech stack built around Transcribe, Translate, and Polly \cite{google_cloud_translate,microsoft_translator,aws_transcribe,aws_translate,aws_polly}. These platforms demonstrate that speech translation has practical value, but they also embody architectural and operational assumptions that are difficult to inspect, adapt, or optimize for smaller actors, research prototypes, and open-source deployment scenarios. This creates a gap between what can be purchased as a managed capability and what can be studied as a transparent, reproducible, and extensible distributed system.

This thesis addresses that gap by investigating an event-driven microservice architecture for real-time speech-to-speech translation. Instead of treating translation as a single opaque model endpoint, the project frames the system as a pipeline of cooperating services for ingress, speech recognition, translation, speech synthesis, orchestration, and evidence collection. That framing is relevant because architectural choices about decoupling, message flow, observability, and artifact handling shape whether the system can achieve low latency, scale under load, and remain robust when components fail \cite{newman2019monolith,richardson2018microservices,bass2021softwarearchitectureinpractice}.

\section{Problem Statement}
\label{sec:problem}

The literature on speech translation is rich in model innovation, but comparatively thinner on deployable system architecture. Research has produced strong advances in neural machine translation, speech recognition, and direct speech-to-speech translation, including transformer-based sequence modeling, large-scale weakly supervised ASR, speech-to-text toolchains, and recent direct translation pipelines \cite{bahdanau2014neural,vaswani2017attention,radford2022whisper,wang2020fairseq,inaguma2023unity,sarim2024direct,zhang2024direct}. However, these advances do not by themselves explain how multiple AI components should be composed into a distributed platform that must satisfy tail-latency constraints, preserve evidence integrity, and remain maintainable over time.

This gap becomes more pronounced when the system must be evaluated as an architecture rather than a benchmarked model. A real-time speech translation platform must coordinate data flow across ASR, MT, and TTS stages; isolate failures; make queueing and service boundaries observable; and support reproducible measurement of latency and robustness under realistic load \cite{kleppmann2017designing,hohpe2004enterprise,bass2023engineeringai,xu2025scalability}. Existing work rarely treats these system-level concerns as the primary object of study. The central problem addressed in this thesis is therefore not only how to translate speech, but how to design and justify an event-driven architecture that makes such translation low-latency, scalable, robust, and empirically assessable.

\section{Research Questions}
\label{sec:research-questions}

The thesis is organized around one main research question and five supporting sub-questions.

\subsection{Main Research Question}

How can an event-driven microservice architecture be designed and implemented to deliver robust and scalable real-time speech-to-speech translation with measurable performance optimization?

\subsection{Sub-questions}

\begin{enumerate}
\item What architectural principles and design patterns support an efficient real-time speech translation pipeline?
\item What factors affect the system's latency and ability to scale?
\item How can these factors be measured and optimized in practice?
\item How can the system be made robust and available even if individual components fail?
\item How can the system's performance and quality be evaluated and possibly compared to existing solutions?
\end{enumerate}

These questions are intentionally distributed across the thesis rather than answered inside the introduction alone. The current chapter frames the problem, scope, and evaluation logic. The background chapter establishes the relevant literature and design vocabulary. Later chapters address architectural response, implementation realization, and empirical evaluation in sequence.


\section{Success Criteria}
\label{sec:succes-criteria}

The thesis evaluates success across architectural, empirical, and methodological dimensions.

\subsection{Real-Time Success}

For this project, \emph{real-time success} is defined by the system's ability to process speech input and return translated output with an end-to-end latency target of P99 below two seconds under the scoped MVP conditions. This criterion does not claim universal real-time performance across all workloads; instead, it provides a concrete architectural driver that shapes service decomposition, observability requirements, and the later evaluation strategy. In practical terms, low median latency is insufficient if tail latency remains unstable, because the user experience of conversation is dominated by worst-case delays rather than best-case responses.

\subsection{Scientific Rigor and Evidence Handling}

The thesis also treats \emph{scientific rigor} as a success criterion. Performance and robustness claims must be grounded in reproducible evidence generated through the project's uv-first validation workflow, so that the later implementation and evaluation chapters can refer to repeatable procedures rather than ad hoc one-off measurements. At the same time, manuscript-facing evidence must remain privacy-safe. This means that shareable evidence is restricted to metadata and aggregated observations, such as percentile latency summaries, correlation identifiers, event types, and bounded artifact sizes, while raw transcripts, raw translated text, raw audio, credentials, and full bearer-style URLs are excluded from the manuscript narrative and supporting handoff materials. In particular, latency discussion in this thesis relies on aggregated metrics such as P50, P95, and P99 rather than raw per-request timestamps.

\subsection{Evaluation Orientation}

Success is finally assessed through the combination of latency, scalability, robustness, and quality evaluation. The introduction therefore frames the thesis as a systems investigation in which translation quality matters, but is interpreted together with architectural consequences such as service-boundary overhead, warm-versus-cold execution behavior, and the observability needed to explain measured outcomes. The later evaluation chapter is responsible for validating these criteria empirically; the introduction only defines the standard against which the system will be judged.

\section{Research Contribution and Scope}
\label{sec:contribution-scope}

\subsection{Research Contribution}

The contribution of the thesis is a defensible reference architecture for real-time speech-to-speech translation that connects distributed-systems design choices to measurable performance trade-offs. More specifically, the work contributes: (1) an architectural framing for event-driven composition of ASR, MT, and TTS services; (2) a maintainability-focused discussion of service boundaries, observability, and contract evolution in AI pipelines; and (3) an evaluation-oriented treatment of how latency, scalability, and robustness claims can be supported with reproducible and privacy-safe evidence \cite{bass2021softwarearchitectureinpractice,bass2023engineeringai,xu2025scalability}. The contribution is therefore not a new translation model, but a structured architectural and empirical argument for how an open, inspectable platform can be designed and assessed.

\subsection{Scope Boundaries}

The project includes the design and implementation of an event-driven platform built from existing AI components, together with performance, scalability, and robustness evaluation under the scoped MVP conditions. It does not include training new foundation models, exhaustive legal or organizational analysis, full production hardening, or a general-purpose comparison against every commercial alternative. Security is treated here as an evidence-governance and privacy-by-design constraint for manuscript-facing artifacts, while deeper runtime security engineering remains the concern of later implementation and deployment work.

\section{Thesis Structure and Question Traceability}
\label{sec:thesis-structure-traceability}

The remainder of the thesis is organized to answer the research questions progressively. Chapter~\ref{chap:background} establishes the literature context for speech translation and event-driven systems. It also introduces AI pipeline maintainability and distributed-systems trade-offs. The Analysis and Design chapter translates that context into architectural drivers, quality attributes, service decomposition, and data-governance choices. The Implementation chapter explains how the platform was realized in practice, including orchestration, artifact handling, and robustness mechanisms. The Evaluation chapter measures the system against the success criteria defined above. The Conclusion then synthesizes the final answer to the main research question.

This structure also makes the sub-question ownership explicit. Sub-question 1 is motivated in the Background chapter and answered architecturally in Analysis and Design. Sub-questions 2 and 3 are framed here first. They are then operationalized through design and implementation and resolved empirically in Evaluation. Sub-question 4 is addressed through implementation and robustness evidence. Sub-question 5 depends on the evaluation and conclusion chapters. The introduction therefore functions as a scope-locking and traceability chapter, ensuring that later arguments remain tied to the original research objective.

